\section{Introduction}
Ground robots usually take their position from a global navigation satellite system (GNSS). When that signal is lost, dead reckoning from inertial and visual sensors has to carry the robot until it returns. The loss is not hypothetical: aircraft operators reported 37,770 events of GPS signal loss in the Middle East region in 2023, 92.96 per 1,000 flights against 29.83 worldwide, in data presented to ICAO \cite{icao_wp69}; these are reported events, not confirmed jamming. Aviation and maritime authorities describe jamming and interference in the Gulf \cite{easa_sib, ukmto, mscio}, and the United Arab Emirates reported an increase in interference cases in 2025 \cite{uae_icao}.

A visual inertial fallback has two weaknesses in a desert. Global heading is unobservable from relative measurements \cite{hesch, vinsmono}, and the camera itself degrades in sand, glare and dust \cite{mer, madmax}. Desert ants solve a similar problem with a sky polarisation compass and a stride integrator \cite{ronacher, wittlinger, wehner03}, and robots built on that idea have homed over metres \cite{lambrinos, antbot}. We take from the ant one transferable idea, an absolute heading from the sky combined with a relative motion estimate, and ask when it helps a stereo visual inertial filter on a desert rover, and how the filter should weight it.

This paper reports the third study of a project whose first two studies used a surrogate front end (Study~1) and rendered desert scenes with a single scattering sky (Study~2). Study~2 found that an ungated sky compass is harmful around solar noon, that weighting by the measured degree of polarisation (DoP) or by innovations helps, and that all filters were overconfident. A red team review of Study~2 found that this rested on a sky without multiple scattering, a compass without bias or tilt, a weighting scheme matched to its own simulator, sky conditions independent of the image degradation, and a reused test set. Study~3 rebuilds the evidence so that each of these can fail.

\textbf{Contributions.} (i) A vector Monte Carlo radiative transfer sky with climatological aerosol, checked against measured desert DoP values \cite{horvath} and against one day of open full sky polarisation images \cite{poughon24}; it brackets the noon blind window of a zenith compass in Abu Dhabi at \sWinAnch{} to \sWinRT{} minutes on 21 June and places such windows at all latitudes within about \sLatMax{} degrees of the equator (Section~\ref{sec:sky}). (ii) A Cram\'er Rao analysis showing that no sky polarisation geometry removes the window when the sun is at the zenith, while a ring of off zenith points shortens it to \sWinRing{} minutes or less. (iii) A compass error model with bias, exact tilt geometry and dust coupled to the image visibility, under which the filter no longer knows the true compass noise. (iv) Nineteen systems, including new robust, two mode, bias state and learned weighting schemes, wheel odometry, a sun sensor and a magnetometer, compared under a protocol frozen before a confirmation run on \sNconf{} trajectories from three environments never used before. (v) A simulated homing study and a drift decomposition. All scenes are RENDERED, all compass data SIMULATED; the only MEASURED inputs are the sky polarisation checks and the aerosol climatology.
